\documentclass{article}
\usepackage{assignment_preamble}

\title{Homework 3}
\author{Ravi Kini}
\date{December 4, 2024}

\begin{document}

\maketitle

\problem
According to the Revised Hierarchical Model, bilnguals learn and process words in their second language by forming a lexical link to their L1, which has a stronger conceptual link (lexicons and grammars) to the things they want to express, and so think in L1, then translate their thoughts to L2.

\clearpage

\problem
We can identify that L1 is being used even when speaking in L2 due to the principle of maximizing common ground. For example, speaker of an L1 that uses pronouns, like English, may use pronouns even when speaking pro-drop L2 languages like Spanish.

\clearpage

\problem
Cognates are words in different languages that are spelled similarly; for these bilinguals exhibit faster processing.

\clearpage

\problem
Even for tasks that do not require information about the gender of nouns like speaking in a language that does not make that distinction, information about gender is still accessed, as can be seen from speakers using masculine adjectives for masculine nouns and vice versa. For example, Boroditsky et al (2003) showed that Spanish and German speakers and Basetti 2007 showed that Italian-German bilinguals and Italian monolinguals associated different words with inanimate objects based on their grammatical gender in their L1.


\clearpage

\problem
In lexical fields with partial overlap, the more frequently used words wins out. For example, French-Dutch bilingual speakers use the French word for bottle.

\clearpage

\problem
Priming is when usage of a lingistic feature in one language leads to the its production in another language. 

\clearpage

\problem
When a bilingual speaker exhibits coordinate unbalanced bilingualism, having a lower proficiency in the L2 than in the L1, this results in some overlap in the areas used for storage and processing, with the L2 being activated in more areas in the right hemisphere and the activation being more diffused. We can detect this by measuring the activation of neurons in areas of the brain related to language while a bilingual speaker uses their L1 versus when they use their L2.

\clearpage

\problem
More code switching is likely to occur when a bilingual speaker is experiencing a lot of emotion, due to different languages having different ways of expressing emotions, usually using the L1 to express stronger emotions, and usually for younger or less proficient speakers. Less code switching is likely to occur as a speaker ages and gains more proficiency, as the speaker becomes more detached from the language.

\clearpage

\problem
An example of a less prestigious variant influencing a more prestigious variant is the transfer of evidentiality into Andean Spanish from Quechua and Aymara as a result of the numbber of bilingual speakers using evidentiality in their L1 language. On the other hand, in India, speakers of lower social classes seek to imitate the speech of higher classes, resulting in an evolution in their speech patterns over time.


\clearpage

\problem
In multilingual societies, acculturation in the process of adapting to the values of a community. For exmaple, immigrants to the United States acculturate by developing greater English proficiency, as many Americans are monolingual, and adhering to American social norms.

\clearpage

\problem
A lack of motivation can result in a failure to acquire the L2 language in a multilingual society; for example, L1 English speakers in Quebec failed to progress past a certain level of proficiency in French. Another reason could be attitudes towards the speakers by the dominant social group. For example, L1 Finnish students in Sweden did worse in their L2 language than those in Australia as they were viewed negatively as a minority group.

\clearpage

\problem
Social motivation impacts the level of proficiency one can gain in a language. For example, English L1 speakesr in Quebec showed limits to how much French proficiency they gained, even after exposure in early childhood, due to being motivated to not learn a second language.

\clearpage

\problem
Language is used to create identities in multilingual communities. For example, the Spanish conquerors of Mexico differentiated themselves from the local Indian population by refusing to teach them Castilian Spanish despite gaining proficiency in the indigenous language.

\clearpage

\problem
Usage of a L2 language by biilnguals has been shown to to result in more rational judgements in comparison to the L1 due to the increased emotional distance bilinguals have towards their L2 language.

\end{document}